\begin{table}[htbp]
\centering
\caption{Lee (2009) Bounds: Sample Selection Correction}
\label{tab:lee_bounds}
\begin{adjustbox}{max width=\textwidth}
\begin{threeparttable}
\small
\begin{tabular}{lcccc}
\toprule
 & \multicolumn{2}{c}{Log prices} & \multicolumn{2}{c}{Distance (km)} \\
\cmidrule(lr){2-3} \cmidrule(lr){4-5}
 & Lower bound & Upper bound & Lower bound & Upper bound \\
\midrule
$g65 \times Pre$ & -0.0460*** & -0.1338*** & 27.0609*** & -19.0176*** \\
 & (0.0131) & (0.0114) & (2.3804) & (2.3672) \\
\midrule
Observations & 625,816 & 625,814 & 625,814 & 625,814 \\
R-squared & 0.1886 & 0.1951 & 0.0015 & 0.0017 \\
Trimming proportion & \multicolumn{4}{c}{0.0882} \\
Item FE & YES & YES & YES & YES \\
\bottomrule
\end{tabular}
\begin{tablenotes}
\small
\item \textit{Notes:} Lee (2009) bounds correct for sample selection arising from
treatment effects on completion rates. 18-month window.
Lower/upper bounds obtained by trimming the outcome distribution in the excess-selected cell.
DiD in completion rate: -0.0882.
Standard errors clustered at the item level.
*** \textit{p}$<$0.01, ** \textit{p}$<$0.05, * \textit{p}$<$0.1.
\end{tablenotes}
\end{threeparttable}
\end{adjustbox}
\end{table}
